% ============================================================
% sec/03baseline_literatura.tex — VERSIÓN INTEGRADA, TRES EJES
%
% OBLIGATORIO: un número objetivo con su fuente, declaración
% explícita de qué se busca igualar o superar, y las diferencias
% de protocolo que impiden una comparación justa.
% ============================================================
\section{Baseline de Literatura}
\label{sec:baseline}

\subsection{Ausencia de resultado publicado sobre este corpus}
\label{subsec:ausencialiteratura}

No existe resultado publicado sobre el corpus objetivo, de modo que el punto
de referencia se toma del trabajo análogo cuya tarea, unidad de observación y
métrica son las más cercanas a las nuestras.

\razon{Esta declaración debe abrir la sección y no ocultarse. El enunciado
prevé explícitamente la sustitución cuando el conjunto carece de literatura
previa, de modo que declararla constituye cumplir el requisito y no eludirlo.
La ausencia se recoge además como brecha en la
Sección~\ref{subsec:brechas}.}

\subsection{Origen y protocolo del punto de referencia}
\label{subsec:origenbaseline}

\begin{quote}
\textbf{Objetivo de referencia:} área bajo la curva ROC en la banda
$0.78$--$0.80$, proveniente del AAIA'17 Data Mining Challenge sobre
\emph{Hearthstone} \cite{janusz_hearthstone_2017}.

\textbf{Tarea:} predecir el ganador a partir de un único estado de juego.
\textbf{Unidad de observación:} un estado de juego. \textbf{Métrica:} área
bajo la curva ROC.

\textbf{Resultado oficial de referencia:} $0.7846$, obtenido por una red
totalmente conectada con dos capas ocultas sobre la representación tabular,
que quedó en el puesto 94. \textbf{Mejor solución enviada:} $0.8019$, un
ensamble construido sobre redes convolucionales.

\textbf{Protocolo:} 3\,250\,000 estados de entrenamiento provenientes de
aproximadamente 65\,000 partidas frente a 750\,000 estados de prueba
provenientes de más de 180\,000 partidas distintas, con 9 mazos en
entrenamiento y 27 en prueba, de modo que el conjunto de prueba contiene
cartas nunca vistas en entrenamiento; puntuaciones preliminares sobre un
5\,\% aleatorio fijo del conjunto de prueba y puntuaciones finales sobre el
resto.
\end{quote}

\subsection{Qué pretende hacer este trabajo con esa cifra}
\label{subsec:objetivobaseline}

El objetivo operativo es que el modelo clásico más fuerte alcance esa banda
de $0.78$--$0.80$ bajo nuestro propio protocolo primario en la ventana
tardía de la partida, y establecer en qué punto de la partida se alcanza por
primera vez esa banda.

No se afirma que superaremos $0.8019$, porque las dos cifras no están
medidas sobre el mismo objeto. Este trabajo declara el rango como referencia
de orden de magnitud y no como umbral obligatorio.

\razon{Declarar el número como cifra a batir cuando procede de otro juego
constituiría una comparación no legítima. Declararlo como orden de magnitud
cumple el requisito de la sección y respeta a la vez el criterio de
comparabilidad.}

\subsection{Comparabilidad del protocolo}
\label{subsec:comparabilidad}

La comparación se considera \textbf{parcialmente comparable}. Coinciden la
naturaleza binaria del desenlace, el uso de información intra-partida y la
evaluación mediante área bajo la curva. No coinciden el juego, las reglas, la
representación del estado ni la población que genera las partidas.

\noindent\textbf{Diferencias que impiden una comparación justa.}
\begin{itemize}
  \item \textbf{Juego distinto.} \emph{Hearthstone} y el \emph{Pok\'emon TCG}
        difieren en sistemas de recursos, condiciones de victoria y
        estructura de tablero, de modo que el contenido informativo de un
        estado no es la misma magnitud.
  \item \textbf{Población de agentes distinta.} Sus estados provienen de
        agentes que actúan al azar; los nuestros, de agentes que compiten y
        son seleccionados por puntuación. Su propio experimento de
        transferencia muestra que esto importa, y cuesta aproximadamente
        $0.04$ de área bajo la curva en la dirección de agentes aleatorios a
        agentes fuertes \cite{janusz_hearthstone_2017}.
  \item \textbf{Partición distinta.} Su separación es por partida y por mazo
        entre dos conjuntos fijos; la nuestra es una partición agrupada por
        episodio, con protocolos adicionales temporal y de reserva de mazos.
  \item \textbf{Prevalencia de clases distinta.} Su conjunto de entrenamiento
        tiene $50.46\%$ de victorias y el de prueba $57.27\%$, de modo que
        las dos mitades de su propio banco de pruebas no provienen de la
        misma distribución \cite{janusz_hearthstone_2017}.
  \item \textbf{Descripción del estado distinta.} Su representación tabular
        fue suministrada por los organizadores, y las mejores participaciones
        reconstruyeron características a partir del JSON crudo; nuestro
        conjunto de características lo construimos nosotros y está
        restringido por una lista negra de fuga.
  \item \textbf{Conjunto de métricas distinto.} Aunque ambos reportan área
        bajo la curva, este trabajo añade pérdida logarítmica, puntaje de
        Brier y curvas de fiabilidad, de modo que la evaluación no se reduce
        a la misma métrica.
\end{itemize}

\razon{Esta enumeración es la exigencia central del criterio de baseline.
Hodge \textit{et al.} ofrecen el precedente: publican una tabla comparativa
de dieciséis trabajos y advierten en nota al pie que la variación en tamaño,
composición y versión del juego obliga a tratar con cautela la comparación de
las exactitudes reportadas \cite{hodge_winprediction_2021}.}

\noindent\textbf{Interpretación del número.} El intervalo de $0.7846$ a
$0.8019$ se emplea como referencia externa del nivel de discriminación
alcanzado en una tarea análoga. No constituye una cifra que este proyecto
deba superar para considerarse válido: un resultado inferior puede ser
metodológicamente sólido si se obtiene bajo un protocolo más estricto, y uno
superior no probaría de forma automática un mejor desempeño, dadas las
diferencias enumeradas.

\subsection{Anclajes secundarios de orden de magnitud}
\label{subsec:referenciasadicionales}

Se dispone de un segundo anclaje, más débil. SC2EGSet reporta $0.9071$ de
acierto en predicción de victoria con dos jugadores y $0.8488$ con un solo
jugador \cite{bialecki_sc2egset_2023}. Se registra aquí por completitud y no
se adopta como objetivo: la métrica es acierto y no área bajo la curva, la
unidad es una partida completa y no un instante, las características son
promedios sobre la partida, y la cuestión de la agrupación planteada en la
Sección~\ref{subsec:comparativo} queda sin resolver.

Hodge \textit{et al.} reportan exactitudes sobre datos mixtos y profesionales
de un videojuego de arena multijugador, evaluando en el punto medio de la
partida \cite{hodge_winprediction_2021}. Xenopoulos \textit{et al.} reportan
valores de pérdida logarítmica y de error de calibración esperado para
modelos evaluados dentro y fuera de su entorno de entrenamiento
\cite{xenopoulos_winprob_2022}.

\pc{Consignar los valores exactos de la Tabla III de Hodge \textit{et al.} y
de la Tabla 1 de Xenopoulos \textit{et al.} tras verificación directa en el
artículo.}

\razon{Estas referencias se incluyen con su métrica original y sin
convertirla. Traducir una exactitud a un área bajo la curva, o una pérdida
logarítmica a una exactitud, produciría cifras que ninguno de los artículos
reporta.}

\decidir{Queda pendiente definir si el trabajo declara además un umbral
propio de éxito, independiente de la literatura, formulado sobre los modelos
de referencia internos descritos en la Sección~\ref{subsec:planbaselines}. Un
umbral de esa naturaleza sería comparable de forma legítima, a diferencia del
número de literatura. La decisión afecta la redacción de la discusión final y
conviene tomarla antes de la Etapa 2.}